% !TeX root = ../papers/nc-p3.tex
\section{片側欠色曲線の無印同型と二つの周期1点}
\label{sec:single-missing-curve-isomorphism}
\begin{theorem}\label{thm:single-missing-curve-isomorphism}
定理\ref{thm:klein-single-missing-windows}の全ての窓を、四元群作用、二組の文字、巡回表示の印を忘れ、
添字保持同型で分類すると、同型類は $X\in k^*$ と一対一に対応する。
基準の始点では $X=(t+t^{-1})^2$、一位置移動は $X\mapsto16/X$、反転は
前節の基準表示を元の類に戻す。最小周期は $X=4,-4$ の二点で1、その他の全点で2である。
\end{theorem}
\begin{proof}[点対応から二直線を復元する]
前節の二組に対応する点空間を $A,B$ と書き、各々二次元とする。次位置では $A',B'$ とする。
内部関係が定める相手への線型写像を $f:A\to A',g:B\to B'$ とする。
例えば内部形式 $a^t D a'=0$ に対し $f=JD^t$、$J$ は二行の標準交代行列である。
各射影線の外の点 $[a:b]$ について、その相手は $[sfa:t'gb]$ と書ける。
混合関係は $s C(a\otimes b)+t' D'(a\otimes b)=0$ になる。
$C$ は可逆で、$M=C^{-1}D'$ と置く。
ここの $a,b$ の両方が非零のとき $a\otimes b$ は $M$ の固有ベクトルでなければならず、
逆にそのとき $s,t'$ の非零な比が一意に決まる。一方 $a,b$ の片方が零の点は全て相手を持つ。

係数表から四位置で $M^2=kI$, $k\ne0$, $\operatorname{tr} M=0$ となる。
二つの固有空間の射影直線をSegre曲面 $\mathbb P A\times\mathbb P B$ と交わらせる。
第一成分 $a=(1,z)$ が交点の端点になる条件は、$P=a\otimes B$ の二列を取ると
$\det[P,MP]=0$ である。実際 $P\cap MP$ が非零なら $M^2=kI$ により交叉自体が $M$ 不変で、
固有ベクトルを含む。第二成分でも同様である。各位置でこの二つの二元四次式は
$az^4+bz^2+c$ ($ac\ne0$) の形で、$H=t+t^{-1}$ として
\[
       b^2/(4ac)=(H+1)^2,\quad (H-1)^2
\]
をそれぞれ与える。従って固有直線は曲面に含まれない。
各四次式が重根を持つ条件は $(H\pm1)^2=1$ で、$H\ne0$ のため少なくとも一方の四次式は
四つの相異なる端点を与える。他側に重根がある場合にもその重複度は2である。
よって常に相異なる四つの分解可能固有直線、従って射影的には四つの点対 $([a],[b])$ を得る。
第一射影の被約支持は $\mathbb P A,\mathbb P B$ と各点対を結ぶ四本の直線の和である。
二つの基準直線は他の四本と全て交わり、各結合直線の方は二つの基準直線と、
高々もう一つの結合直線と交わる。従って無印の辺から基準直線の無順序対が全点で復元される。
\end{proof}

\begin{proof}[一辺の数値不変量と二窓の比較]
$A\otimes B$ 上に二行二列行列の行列式形式を取り、その随伴を $M^\dagger$ とする。
基準直線対の同型はこの形式の相似変換になる。二つの組を交換すると、
混合式の二項が入れ替わるため、因子の交換を除けば $M$ は $M^{-1}$ になる。
$M^{-1}=M/k$ であり、内部写像のスカラー選択も $M$ をスカラー倍するだけである。
よって
\[
            I=\operatorname{tr}\big((M^\dagger M/k)^2\big)
\]
は無印の辺の不変量である。全四位置で順に
$4-32/X,4-2X,4-32/X,4-2X$ を得る。基準始点の同型では $X$ の一致が必要である。
逆に $X$ が同じならパラメータは $t,-t,t^{-1},-t^{-1}$ のいずれかである。
$-t$ と $t^{-1}$ への同型は四位置のラベル置換と倍率で実現する。
具体的には置換を $0,1,2,3$ の像の列として、$-t$ に対し位置0,1で $2301$、
位置2,3で $3210$、$t^{-1}$ に対し位置0,1で $0132$、位置2,3で $1023$ を取る。
残る非零成分比は四位置各四個の倍率の指標系で、その全整数左核積は1である。
二位置移動は四位置全てで $2310$、四文字の逆順は前二位置 $2310$、後二位置 $3201$ により
同じ指標検査を満たす。従って反転と二移動はパラメータ類を増やさない。

一移動の十分性には混合四行列そのものを比較する。各次位置の $A,B$ を前の位置の $f,g$ で
引き戻し、最初の三辺の $M$ を $M_0,N_1,N_2$ と同じ四空間で読む。
$N_2=-4t^3M_0/(1+t^2)$ である。$i^2=-1$, $s=i(t-i)/(t+i)$, $L=1+t^2$ として
\[
 P=\begin{pmatrix}1&-it/L\\-4it/L&4t^2/L^2\end{pmatrix},
 \qquad Q=\begin{pmatrix}1&-1/t\\-s&-s/t\end{pmatrix},\qquad C=P\otimes Q .
\]
二つの可逆行列は全許容値で
\[
 C M_0(t) C^{-1}\sim N_1(s),\qquad
 C N_1(t) C^{-1}\sim M_0(s)\sim N_2(s)
\]
を満たす。$\sim$ は非零スカラー倍を表す。この比較で二窓の同型が従う理由を述べる。
内部形式を両側の $f,g$ で同定した後、混合関係四空間は二つの混合四空間の間の
可逆写像のグラフであり、その行を左から $C_{\rm left}^{-1}$ で揃えたものが上で用いた $M$ である
（$C_{\rm left}$ は混合式の第一係数行列）。従って $P,Q$ の基底変更でグラフの写像がスカラー倍まで一致する。
二辺のスカラーは、三つの位置で二組の基底に掛ける相対倍率の隣接比で順に吸収できる。
内部関係は各々の直線として保たれる。これで二つの連続窓が同型となり、
系\ref{cor:two-window-extension}で添字保持同型へ延びる。
$(s+1/s)^2=16/X$ なので全一移動が基準曲線に戻り、固定点は $X^2=16$ の二点である。
最後に $X\ne0$ なら $t+t^{-1}=\sqrt X$ は許容根を持つのでパラメータの像は $k^*$ 全体である。
\path{computations/verify_single_missing_isomorphism.py} は八個の四次式、四位置の不変量、
置換後の整数左核および二つの基底変更恒等式を有理関数体で検算する。
\end{proof}
